\section{Indispensable Spiritual Qualities}

\textbf{Frithjof Schuon} concludes the essay \emph{Modes of Spiritual Realization} with a succinct description of the ``qualities that are indispensable for spirituality.'' In preparation, he criticizes both \textbf{moralism} and \textbf{rationalism}.


\paragraph{Moralism}


Moralists, or humanitarians, are ``indignant over the indifference which the saints --- both western and eastern --- show towards human misery,'' but do not understand the reasons why.

\begin{itemize}


\item Many of the miseries in the traditional world are lesser evils since some calamities are inevitable. The moralistic and humanitarian modern men, in their frenzied activity, fail to grasp that some events simply cannot be avoided. The artificial suppression of such events results in worse unintended consequences.

\item The saints are indifferent to contingent events because they desire to deal with evil at its root by returning to the Source, not by dissipating energy in illusory efforts.
\end{itemize}


Contingent evils may be ameliorated, but Evil itself is unavoidable. Christ healed the sick, but did not eliminate sickness itself.

\paragraph{Rationalism}


Newcomers to Tradition assume it is just another philosophical position to be debated, supported, or refuted in thought. They cannot conceive of anything deeper and thus completely misunderstand Tradition. Schuon explains:

\begin{quotex}
In metaphysical knowledge, reasoning can play no other role than that of occasional cause of intellection; the latter intervenes in a sudden, not continuous or progressive, way as soon as the mental operation, conditioned in turn by an intellectual intuition, possesses the quality or perfection that makes of it an effective symbol.
\end{quotex}


That is, one ``gets it'' in a flash; one is not convinced by a proof or the accumulation of evidence. Rationalism, on the contrary, looks for Truth as a formulation of thoughts; it rejects a priori the possibility of a higher knowledge or gnosis beyond thought. Yet any mental formulation of the world cannot be the world itself, just as the map is not the territory.

The Rationalist tries to find a standpoint outside the world, from which to formulate a picture of the world, so he necessarily leaves himself out of his own picture! But, Pure Truth is identical with that which is, and thereby includes the thinker.

\paragraph{Qualities}


There are three qualities the Schuon considers indispensable for spirituality.

\begin{description}

\item[Objectivity]

This is a perfectly disinterested attitude of the intelligences that is free from ambition and prejudice, hence accompanied by serenity.
\item[Magnanimity]

This is nobility or the capacity of the soul to rise above all things that are petty and mean.
\item[Simplicity]

Man is freed from every unconscious complex or compulsion stemming from self-love. He is without artifice, pretension, ostentation, dissimulation, and pride.
\end{description}


Obviously, these qualities are irrelevant to the thinker, the philosopher, the politician, the activist, the scientist.

\flrightit{Posted on 2011-02-03 by Cologero}

\begin{center}* * *\end{center}

\begin{footnotesize}
\begin{sffamily}
\texttt{Jason on 2011-08-12 at 08:29 said:}

I normally prefer Guenon to Schuon, but he's spot-on here.

I haven't read enough of him to know the answer, but what of the Christian doctrine of the second coming, that sickness and Evil itself will be eliminated at a yet-future time?

I am not convinced that the eschaton as it is understood in Christian theology actually means the end of struggle. It is certainly not the end of work, if the imagery in Revelation 22 of the kings bringing the ``glory of the nations'' into the Heavenly City is to be understood (as I understand it) as human creative activity becoming a part of the City.

Perhaps it might be more pertinent to say that ``Evil'' is not this metaphysical, moral wrongness we in the West often ascribe to it, but more of a struggling -- that which we overcome to achieve a meaningful life.

\hfill

\end{sffamily}
\end{footnotesize}

